\section{Conclusion}
\label{sec:conclusion}

We measured six quantitative claims that circulate about CNN--transformer
hybrids and state space vision backbones. The cross-over in computational
cost is where the analysis says it is. The size of the advantage past it is
not: counted end to end on whole models it is a factor of two rather than an
order of magnitude. The reported throughput and memory advantages did not
reproduce on our hardware, and at batch size one the memory ordering is
reversed. The receptive field is anisotropic along the scan direction, but
only on a metric sensitive to the core of the field, and the effect is
present before any training. The freedom from softmax normalisation does not
put more attribution on large objects; the state space model places the least
of the three at every object size we measured. And the accuracy gap does not
have to be explained by a decomposition borrowed from one architecture: a
factorial design says the structural prior dominates, that the choice of
token-mixing operator matters far less at this sequence length, and that the
two interact in a direction we can sign but not size.

Four of the six do not hold as stated, and in three of those the argument
behind the claim is not what fails. What fails is the identification of the
argument with the number reported beside it---a kernel counted as a block, a
field's skirt measured for its core, an absent normalisation read as a
present capacity. The useful output of an audit like this is not a list of
corrections but the axis on which each claim is true, stated so that the next
comparison can quote it without repeating the substitution.

What remains open is stated in Section~\ref{sec:limitations}, and the two
items we would run next are the flat cells at a sequence length past the
cross-over, and a three-factor design that separates hierarchical
downsampling from convolutional locality.
